\section{Rendezvous / Bonjour}
\label{sect:bonjour}
\marginpar{\includegraphics[width=1.25cm,height=1.25cm]{../../images/net_bonjour.eps}}

\emph{Bonjour}, formerly known as \emph{Rendezvous}, enables automatic discovery of computers, devices, and services on IP networks. It uses industry standard IP protocols to allow devices to automatically discover each other without the need to enter IP addresses or configure DNS servers. It is a technology developed by Apple that is submitted to the IETF as part of an ongoing standards-creation process. The technology is more generally known as zero-conf networking.

There are opensource \emph{Java} libraries available that allow to use this technology in any \emph{Java} application. The library needs a native library that is available for all major platforms. 

DNS service discovery is a way of using standard DNS programming interfaces, servers, and packet formats to browse the network for services. It is compatible with, but not dependent on, multicast DNS. Multicast DNS is a way of using familiar DNS programming interfaces, packet formats and operating semantics, in a small network where no conventional DNS server has been installed.


\subsection{Prototype}
The technology is best described by a simple prototype. We've created a simple echo server and echo client that show off the basic features of this technology.

\subsubsection{Echo Server}
The echo server is implemented in the class \texttt{ch.iserver.ace.net.bonjour.EchoServer}. It is very similar to normal echo server implementation. First a \texttt{ServerSocket} is created listening on a particular socket. Then the service is registered with DNS service discovery system represented by the \texttt{com.apple.dnssd.DNSSD} class. A listener (\texttt{com.apple.dnssd.RegisterListener}) is used to inform the application when the registration was successful. Once the registration was successful, the echo server starts to service clients as any echo server would do.

\subsubsection{Echo Client}
The echo client (see \texttt{ch.iserver.ace.net.bonjour.EchoClient}) is a simple GUI that allows to send echo requests to echo servers. It browses the network for available echo servers and lists them in a combo box. The user can select one of those echo servers and send an echo request to them.

The method \texttt{browse} of the \texttt{DNSSD} class is used to issue a browse request to the DNS service discovery system. We have to supply a \texttt{com.apple.dnssd.BrowseListener} implementation that gets called back whenever a service of specified type is found.

Whenever the user sends a message to the currently selected echo server by entering a message in the bottom textfield and hitting enter, the service has to be resolved. This happens by calling \texttt{resolve} on the \texttt{DNSSD} class. The application is called back over a \texttt{com.apple.dnssd.ResolveListener}. The method \texttt{serviceResolved} is called on this listener when the service is resolved. This provides us with the service name, host name and port of the service. As a last step we have to get the IP address of the service. This is done by calling \texttt{queryRecord} on the \texttt{DNSSD} class. We get called back by a method call to \texttt{queryAnswered} in the \texttt{com.apple.dnssd.QueryListener} we supplied to the \texttt{queryRecord} call.

At that point we do have enough information to contact the echo server and send the echo request over a plain socket connection. That is one important point to observe. This technology enables us to discover services on the network but does not impose any limitations on the way we communicate from client to server.


\subsection{Analysis}
\emph{Bonjour} technology provides discovery of other services on a LAN. It leaves the communication protocol completely up to the service. So for instance a client can discover a HTTP server and communicate by sending HTTP requests with the server. He can discover a SSH server and then communicate by talking the SSH protocol. So \emph{Bonjour} does not provide any support for communication.

SubEthaEdit, the only commercially available collaborative editor for OS X, uses \emph{Bonjour} to discover other published documents. Once other documents are discovered, SubEthaEdit uses \emph{BEEP} (see section \ref{sect:beep}) using a proprietary protocol to communicate with other instances of the application. This combination could also be a viable alternative for our network layer implementation.

\subsubsection{Positive Points}
\emph{Bonjour} provides a platform and programming language independent way of discovering other services on the network. It is on its way of becoming an IETF standard. This technology will remain available for the years to come.

This would allow other applications programmed in other programming languages to interoperate with \ace. Depending on our requirements, interoperability could be a very important point.

\subsubsection{Negative Points}
\emph{Bonjour} discovery gets you only a host name plus a port. Of course, this information can be used to open a socket, but other technologies provide you simpler ways of communicating with each other (e.g. \emph{RMI}).

The discovery is further limited to the LAN. Although it would be theoretically possible to extend this behavior, it needed fairly deep understanding of the technology.

\subsubsection{Open Points}
Implementing discovery using \emph{Bonjour} would be pretty straightforward. There are no open points to examine.


\subsection{Resources}
\begin{itemize}
 \item \href{http://developer.apple.com/darwin/projects/bonjour/}{Developer site for \emph{Bonjour} technology}
 \item \href{http://www.macdevcenter.com/pub/a/mac/2004/08/31/osx_java.html}{Article about Java \emph{Bonjour}}
 \item \href{http://developer.apple.com/documentation/Java/Reference/DNSServiceDiscovery_JavaRef/index.html}{Java API documentation}
\end{itemize}
